\section{From a coherent answer to a defensible causal decision}
\label{exec17:frontier-start}

Three recent systems clarify different meanings of assistance. A verifier can
check whether an answer follows from the reasoning presented to it; an
orchestrator can choose and run a sequence of causal tools; a literature agent
can extract candidate relations and draw a diagram. These are valuable but
different contributions. A bank considering a retention offer needs an answer
to a further question: under which assumptions would that offer change a
specified outcome for a specified population? We compare the mechanisms of
Recursive Causal Audit, OrcheCause and EpiCausalX against this decision, keeping
reported demonstrations distinct from a proposed banking implementation.
\citep{exec17_frontier_chang,exec17_frontier_yang,exec17_frontier_liu}

\subsection{Process consistency and the scientific oracle}

Chang's Recursive Causal Audit, or RCA, addresses a specific failure: a model
can describe a problem in one way and then emit a contradictory final label,
particularly when challenged by a user. Its judge has no gold labels. It checks
the output schema, consistency among fields, whether the conclusion follows
from the structured reasoning, and whether a user hint dominates contrary
reasoning. Staged prompts demand more explicit variables, causal sketches and
assumptions when the earlier output fails. The paper explicitly recognizes
that a consistently incorrect argument can pass the judge. That limitation is
part of the mechanism's definition, rather than an accidental omission that
another formatting check can repair.\citep{exec17_frontier_chang}

In a retention analysis, the agent might state that adviser discretion creates
confounding and then report the treated-minus-untreated retention difference
as an identified offer effect. A consistency judge could catch the conflict.
It would not necessarily catch a different response that consistently asserts
random assignment throughout, when advisers actually selected customers using
unrecorded risk. The latter response needs evidence about the assignment
process. Agreement between a rationale and a number is a necessary internal
check for a well-formed report, but does not supply that evidence.

The study uses a curated 454-vignette diagnostic suite and standardized pressure
protocols. Its limitations include subjective annotations, prompt dependence,
black-box explanations of model behavior and the absence of a component-level
ablation of RCA. The historical calls run through January 2026; they are not a
ranking of models available at the time of this proposal. There is also a
specific unresolved denominator issue: the statistical-analysis paragraph
describes L1 with 50 cases per class and 100 overall, while Table 7 totals 50 L1,
304 L2 and 100 L3 cases. We therefore do not reconstruct a pooled confidence
interval or repair the denominator by assumption.\citep{exec17_frontier_chang}

The answer-transition metric requires equally careful scope. In the paper's
Table 4, the L2 labels are binary, VALID and FLAWED. For a fixed gold label,
changing an incorrect binary answer necessarily corrects it. Let $p$ denote
initial accuracy, $b$ the probability of becoming wrong conditional on an
initially correct answer, and $c$ the probability of becoming correct
conditional on an initially wrong answer. Then
\begin{equation}
 p_{\mathrm{final}}=p(1-b)+(1-p)c.
 \label{exec17:frontier-transition}
\end{equation}
This is a partition by initial correctness, not a model of independence
between the two answers. Substitution of the three rounded rows in Table 4
agrees with their reported final accuracies to within 0.1 percentage point.
It is therefore unwarranted to call this particular binary table erroneous
merely because a more general answer vocabulary exists elsewhere in the
paper.\citep{exec17_frontier_chang,exec17_frontier_exact}

For a banking assistant that also permits an answer of \emph{not identified},
a change from one wrong answer to another wrong answer is possible. To see
the difference, fix a task with true label $C$ and two wrong labels $W_1,W_2$.
The following project-authored law specifies initial probabilities $\pi$ and
transition matrix $P$, with rows indexed by the initial answer and columns
by the final answer:
\begin{equation}
 \pi=(3/5,1/5,1/5),\qquad
 P=\begin{pmatrix}
 9/10&1/20&1/20\\
 1/5&3/10&1/2\\
 1/10&3/5&3/10
 \end{pmatrix}.
 \label{exec17:frontier-multiclass}
\end{equation}
All rows sum to one. Among initially wrong answers, the probability of any
label change is $7/10$, whereas the probability of correction is only $3/20$.
The bad-flip probability is $1/10$. Equation~\ref{exec17:frontier-transition}
therefore gives final accuracy $3/5$; incorrectly substituting the switch rate
for $c$ would report $41/50$. Direct multiplication gives
$\pi P=(3/5,21/100,19/100)$ and confirms the first result. These are exact
constructed probabilities, not measured model performance.
\citep{exec17_frontier_exact}

The practical implication is to preserve a full answer transition table,
including refusal and tool failure, with a fixed target-specific gold label.
Whether refusal is correct depends on whether the question is identified
under the supplied assumptions. A pressure test should separately expose
unsupported endorsement, inappropriate rejection and appropriate withholding.
Its denominator must include the same assigned questions before and after
review; deleting failed or abstaining requests can change the apparent benefit.

\subsection{Confounding is not a proof of a zero effect}

A vignette in Chang's appendix describes treatment given to more severely ill
patients. It labels the positive causal assertion NO and describes the
association as spurious. The defensible lesson is that the supplied
association does not establish the asserted effect. Confounding by itself
does not establish that the treatment has no effect. A gold-label protocol
must distinguish an unjustified assertion from a false effect assertion.
\citep{exec17_frontier_chang}

Consider the following authored banking construction. Let $A$ indicate a
retention offer, $Y$ indicate retention over a fixed horizon, and $U$ be an
unrecorded customer characteristic. Draw $U$ as Bernoulli with probability
$1/2$. Independent exogenous uniform variables generate the offer and
outcome using the conditional probabilities
\begin{equation}
 \Pr(A=1\mid U=u)=\tfrac14+\tfrac12u,\qquad
 \Pr(Y=1\mid U=u,A=a)=\alpha+\beta u+\tau a.
 \label{exec17:frontier-scm}
\end{equation}
The outcome equation is retained under intervention; only the offer equation
is replaced. All conditional probabilities below lie strictly between zero
and one, so the comparison does not depend on deterministic treatment or
an empty observed treatment stratum.\citep{exec17_frontier_exact}

\begin{center}
\begin{tabular}{lrrrrr}
\toprule
World & $\alpha$ & $\beta$ & $\tau$ & $E[Y(0)]$ & $E[Y(1)]$\\
\midrule
Positive effect & $7/20$ & $1/5$ & $1/10$ & $9/20$ & $11/20$\\
Zero effect & $3/10$ & $2/5$ & $0$ & $1/2$ & $1/2$\\
Negative effect & $1/4$ & $3/5$ & $-1/10$ & $11/20$ & $9/20$\\
\bottomrule
\end{tabular}
\end{center}

Each world has $\Pr(U=1\mid A=1)=3/4$ and
$\Pr(U=1\mid A=0)=1/4$. Consequently the observed conditional means are
\begin{equation}
 E[Y\mid A=0]=\alpha+\beta/4=2/5,\qquad
 E[Y\mid A=1]=\alpha+3\beta/4+\tau=3/5
 \label{exec17:frontier-observed}
\end{equation}
in all three worlds. The complete observed joint law of $(A,Y)$, not only
its two means, is identical: the cells $(0,0),(0,1),(1,0),(1,1)$ have
probabilities $3/10,1/5,1/5,3/10$. Under intervention,
$E[Y(a)]=\alpha+\beta/2+\tau a$, so the causal difference is $\tau$.
Thus the same positive observed difference of $1/5$ coexists with a
positive, zero or negative effect. Independent finite enumeration of the
threshold construction reproduces both the observed and intervention laws.
\citep{exec17_frontier_exact}

For a stipulated USD 100 value per additional retained customer and USD 5
cost per offer, incremental value is $100\tau-5$: USD 5, minusUSD 5 or
minus USD 15 in the respective worlds. These assumed amounts illustrate why
an unjustified zero-effect conclusion and an unjustified positive-effect
conclusion can both mislead a decision. They are not estimates of bank prices
or returns. Even unlimited observations of $(A,Y)$ would leave these three
worlds observationally indistinguishable. The next investment must supply
defensible restrictions or new information, such as an appropriately designed
public or synthetic randomized demonstration before proposing an expensive
bank trial.\citep{exec17_frontier_exact}

\subsection{Text constraints and the limits of graph-score tuning}

OrcheCause offers a complementary mechanism: task orchestration across causal
discovery, structural causal modeling, root-cause analysis and causal-effect
estimation. Its paper describes causal-learn, DoWhy, pgmpy and causalml as
underlying libraries. User text yields hard cause--effect constraints, while
external database text yields soft constraints. Preprocessing includes
imputation and removal of highly correlated features; hyperparameter search
can use BIC when a ground-truth graph is unavailable.
\citep{exec17_frontier_yang}

This distinction between hard and soft text makes a concrete design choice
visible. A hard requirement $A\to Y$ excludes every graph without that edge.
If the sentence actually describes a historical association, a mediated
effect, an intended policy or a different time scale, the true explanation
may already have been removed before a statistical test runs. A soft
constraint can permit disagreement with the data, but its weight is still
a modeling choice; an LLM's linguistic confidence is not automatically a
calibrated probability that the edge is true. In the proposed bank interface,
each relation therefore needs an explicit role: observed association,
mechanistic hypothesis, temporal prohibition, intervention definition or
accepted scientific restriction. The strength of a constraint must follow
its accepted meaning and scope.

The paper's table also argues against a single undifferentiated improvement
score. The following are the reported PC demonstrations; HPO denotes
hyperparameter optimization, DK domain knowledge from the original dataset
papers, and EK external knowledge. Lower structural Hamming distance (SHD)
and higher F1 are favorable under the paper's respective metrics.
\citep{exec17_frontier_yang}

\begin{center}
\begin{tabular}{llrrrr}
\toprule
Dataset & Metric & Baseline & HPO & DK+HPO & DK+EK+HPO\\
\midrule
ALARM & SHD &13&13&10&10\\
      & F1  &.927&.927&.940&.940\\
Sachs & SHD &21&23&14&10\\
      & F1  &.323&.343&.611&.688\\
\bottomrule
\end{tabular}
\end{center}

HPO alone makes Sachs SHD worse, from 21 to 23, while F1 improves. Thus
the paper's general statement of progressive improvement is too broad if
read as improvement of every metric at every stage. Nor does knowledge
drawn from the dataset's own original paper simulate a previously unknown
bank mechanism. The demonstration supports the value of investigating
knowledge-assisted discovery; it does not estimate deployment accuracy,
causal-effect coverage or the incremental value of each component on new
banking questions. Time-series analysis, high-dimensional low-sample
robustness and improved extraction from real documents are future directions
in this paper.\citep{exec17_frontier_yang}

Observed fit cannot by itself solve this last step. The three complete laws
in Equation~\ref{exec17:frontier-observed} give the same likelihood for every
sample of the observed variables, while their intervention values differ.
No criterion computed solely from that observed distribution can select the
correct intervention value without additional restrictions. A graph score
may guide a search inside an assumed family; the product must still expose
the family and its exclusions. Similarly, removing a highly correlated
variable needs a causal-role review: predictive redundancy does not prove
that an adjustment variable or measurement proxy can be discarded.
\citep{exec17_frontier_exact}

For a prototype, the useful adoption is the modular orchestration idea and
the explicit distinction between text sources. Reproducing the reported
implementation requires a separate pinned-code and dependency review.
A first comparison would hold the dataset, target and compute budget fixed,
then compare no text, verified constraints, erroneous constraints and
contradictory documents. The outcome is not merely graph similarity: it
includes whether the effect is identified, whether a wrong constraint
changes the decision, and whether the system discloses the conflict.

\subsection{Literature extraction, diagrams and query-relative roles}

Liu and colleagues' EpiCausalX preprint describes a literature-driven
workflow with natural-language initiation and diagram construction. Its
methods extract relationship terms with regular expressions and infer
directions through heuristics informed by text and epidemiological theory.
The diagram tool checks acyclicity and uses topology to label exposures as
nodes without incoming edges and outcomes as nodes without outgoing edges.
The reported usability exercise includes 12 participants, split into two
groups of 6, and an average task initiation within two dialogue rounds.
These are workflow observations; they are not measurements of causal-effect
accuracy or coverage.\citep{exec17_frontier_liu}

The topology rule illustrates an important boundary for business translation.
In the graph $U\to A$, $U\to Y$, $A\to Y$, the offer $A$ is still the
exposure of interest despite having an incoming edge. If retention $Y$
subsequently affects account value $V$, then $Y$ remains the outcome for a
retention question despite having an outgoing edge. Roles are relative to
the query and intervention, whereas root and sink status are graph
properties. The authored worlds above instantiate the first graph with
explicit positive probabilities. A drawing tool can validate that this
representation is acyclic; it cannot thereby prove the absence of omitted
common causes or identify the effect from observed $(A,Y)$ alone.
\citep{exec17_frontier_exact}

Literature extraction can nevertheless lower the cost of preparing a review.
A useful output is a candidate relation with its original passage, population,
time ordering, measurement definition and study design. An observational
association should remain an association until the proposed causal argument
and its assumptions have been checked. Negation, hypothetical wording,
direction reversal, selection and population mismatch deserve explicit
adversarial examples. The agent can propose a diagram and alternatives, then
pass a well-specified query to a separately reviewed identification routine.
The diagram is an argument to inspect.

The preprint itself places systematic expert review, real-world validation,
larger user studies and comparisons with manual work in its future program.
It also describes open-source release as a future step after initial
validation. Its displayed rights notice and the absence of a reviewed
software release prevent a direct-adoption decision here. The appropriate
proposal is to assess the documented workflow and develop independent
query-typing tests; the paper's usability result does not justify porting an
uninspected implementation or claiming current software availability.
\citep{exec17_frontier_liu}

\subsection{A public-evidence prototype that distinguishes these contributions}

The three mechanisms lead to a staged, falsifiable prototype. Its first
component is a collection of stipulated causal worlds and business questions,
including the equal-observation example, query-relative variable roles,
binary and three-label review tasks, and deliberately false textual
constraints. These fixtures have known answers because their laws are
specified. A second component uses public datasets and documented assumptions
to test tool invocation, provenance and the treatment of missing identification
information. Public observational data need not have known effects: the
expected response may be an identified set, sensitivity analysis or an
explicit inability to decide.

The comparison separates the three interventions. Adding a consistency judge
tests whether contradictory conclusions decrease at a declared review cost.
Adding text constraints tests whether correct information helps and wrong
information is exposed. Adding literature extraction tests whether supported
relations, populations and temporal definitions are preserved. Requests,
failed calls, refusals, changed assumptions and review time must remain in
the evaluation record. For generated mechanisms, replicate mechanisms rather
than treating repeated prompts about the same mechanism as independent
scientific evidence. Predeclared tolerances and held-out problem families
would be needed before estimating an improvement.

The resulting investment case would establish whether these separate
contributions merit a larger experiment. It would preserve the distinction
between coherent reasoning, a diagram, a successful tool call, a better
graph score and a justified banking decision. The attainable next step is
a narrowly specified prototype with known failures and honest uncertainty,
followed by bank-specific validation if that evidence warrants the cost.
\label{exec17:frontier-end}
